\documentclass{beamer}
\beamertemplatenavigationsymbolsempty
\usepackage{tabularx}
\usepackage{tikz}
\usepackage{hyperref}
\hypersetup{colorlinks,linkcolor=,}
\usepackage[type={CC},modifier={by-sa},version={4.0},]{doclicense}
\usetikzlibrary{shapes.geometric, arrows}
\tikzstyle{startstop} = [rectangle, rounded corners, minimum width=3cm, minimum height=1cm,text centered, draw=black, fill=red!30]
\tikzstyle{io} = [trapezium, trapezium left angle=70, trapezium right angle=110, minimum width=3cm, minimum height=1cm, text centered, draw=black, fill=blue!30]
\tikzstyle{process} = [rectangle, minimum width=3cm, minimum height=1cm, text centered, draw=black, fill=orange!30]
\tikzstyle{decision} = [diamond, minimum width=3cm, minimum height=1cm, text centered, draw=black, fill=green!30]
\tikzstyle{arrow} = [thick,->,>=stealth]

\title{\LaTeX \ for Theses}
\subtitle{An Introductory Course}
\author{Dr Thomas Bishop \and Dr Gerard Capes}
\institute{University of Manchester}
\date{17\textsuperscript{th} of November, 2016}
\logo{\includegraphics[width=0.15\paperwidth]{UniOfManchesterLogo.pdf}}

\begin{document}

\begin{frame}
	\titlepage
\end{frame}

\section{Introduction}

\begin{frame} \frametitle{Staff}
	\begin{block}{Tom Bishop}
			\begin{itemize}
				\item Laboratory Technician, Geography
					\item \url{thomasbishop.uk}
					\item \url{github.com/tombishop1}
			\end{itemize}
		\end{block}
	\begin{block}{Gerard Capes}
		\begin{itemize}
			\item Research Software Engineer, IT Services
			\item \url{github.com/gcapes}
		\end{itemize}
	\end{block}
\end{frame}

\begin{frame}
	\frametitle{Course Structure}
	\tableofcontents
\end{frame}

\begin{frame} \frametitle{What We \textbf{Won't} Be Covering}
	\begin{itemize}
		\item How to use \LaTeX \ to typeset other documents (papers, reports, slides, \textit{etc}.),
		\item How to install \LaTeX \ on a Windows/Mac/Unix system,
		\item Every possible option, setting, package or trick,
		\item Image editing or drafting diagrams.
	\end{itemize}
\end{frame}

\section{What Is \LaTeX?}

\begin{frame} \frametitle{What Is \LaTeX?}
	\begin{itemize}
		\item Software for typesetting.
		\item Totally free and available on many platforms.
		\item Popular in publishing technical and/or large documents.
		\item \textbf{A markup language (content and form are separated).}
	\end{itemize}
\end{frame}

\begin{frame} \frametitle{Is it For Me?}
	\begin{columns}
		\column{0.5\textwidth}
			\begin{itemize}
				\item Produces beautiful, professional documents.
				\item There is almost certainly an elegant solution to your typesetting problem.
				\item You can logically work out how to solve an error.
			\end{itemize}
		\column{0.5\textwidth}
			\begin{itemize}
				\item What you see isn’t what you get as you write and edit.
				\item You might have to research how to solve the problem.
				\item You will need to learn how to work with markup language.
			\end{itemize}
	\end{columns}
\end{frame}

\section{How \LaTeX \ Works}

\begin{frame} \frametitle{The Software}	
	\begin{center}
		\begin{tikzpicture}[node distance=0.15\paperwidth]
			\draw [black] (3.5,-1) rectangle (-2,-4.75);
			\node at (2.5,-4.25) {MiK\TeX};
			\node (start) [startstop] 								{Markup Text (from \TeX nicCenter)};
			\node (pro1)  [process, below of=start] 	{\LaTeX};
			\node (pro2)  [process, below of=pro1] 		{\TeX};
			\node (stop)  [startstop, below of=pro2] 	{PDF Output};
			\draw [arrow] (start) -- (pro1);
			\draw [arrow] (pro1) -- (pro2);
			\draw [arrow] (pro2) -- (stop);	
		\end{tikzpicture}
	\end{center}
\end{frame}

\section{The University Thesis Document}

\begin{frame} \frametitle{The University Thesis Document Files}
	\begin{block}{Style Sheet}
		Tells \LaTeX \ how to compose the document.
	\end{block}
	\begin{block}{Thesis Itself}
		The ``backbone'' of the document to which you add chapters.
	\end{block}
\end{frame}

\begin{frame} \frametitle{How do Files Link?}
\begin{center}
		\begin{tikzpicture}[node distance=0.15\paperwidth]
			\node (pro1)  [process] 	{Thesis File};
			\node (pro2) [process, right of=pro1, xshift=2cm] {Chapter 1};
			\node (pro3) [process, below of=pro2] {Chapter 2};
			\node (pro4) [process, right of=pro3, xshift=2cm] {\includegraphics[width=1cm,keepaspectratio]{example-image-a}};
			\node (pro5) [process, below of=pro1] {Bib\TeX \ Database};
			\draw (pro1) -- (pro2);
			\draw (pro1) -- (pro3);
			\draw (pro3) -- (pro4);
			\draw (pro1) -- (pro5);
		\end{tikzpicture}
	\end{center}
\end{frame}

\section{Getting Help and Helping Yourself}

\begin{frame} \frametitle{The Minimum Worked Example}	
	\begin{itemize}
		\item Designed so other users can understand what you are talking about and re-create the problem.
		\item No-one is going to troubleshoot hundreds of lines of code.
		\item Remember you are using a custom stylesheet --- try to provide MWE's in native document classes.
	\end{itemize}
\end{frame}

\begin{frame} \frametitle{Resources}
	\begin{itemize}
		\item \href{http://www.github.com/tombishop1/thesis-latex-manc-coursematerials}{The Course Resources}
		\item \href{http://tex.stackexchange.com}{\LaTeX \ Stack Exchange}
		\item \href{https://wch.github.io/latexsheet/}{\LaTeX \ Cheat Sheet}
		\item \href{https://en.wikibooks.org/wiki/LaTeX}{The \LaTeX \ Wikibook}
		\item \href{https://www.ctan.org/pkg/excel2latex?lang=en}{Excel to \LaTeX \ Converter} and \href{http://www.tablesgenerator.com/}{\LaTeX \ Table Maker}
	\end{itemize}
\end{frame}

\begin{frame} \frametitle{Concluding Remarks}
	\begin{itemize}
		\item You (probably) don't have to decide now.
		\item It takes time to learn, but you'll spend time typesetting in any program, and it's a useful skill.
		\item Discuss feedback, notes, version control and typesetting with your supervisor as soon as you can.
	\end{itemize}
\end{frame}

\begin{frame}
\doclicenseThis
\end{frame}

\end{document}